\chapter{Experimental Setup}
\label{chap:expsetup}

Translating moodboard inputs into effective conditioning signals for 3D generative models presents a fundamental architectural challenge. Whereas the preceding chapter detailed how Imagin3D orchestrates this translation, the present chapter describes how that process is evaluated. It first establishes the overarching evaluation strategy, motivating the hybrid approach that combines objective embedding metrics with subjective human studies. It then formally defines the baseline workflow against which Imagin3D is compared, specifies the composition and structure of the evaluation dataset, and details both the subjective and objective protocols used to measure performance.

\section{Evaluation Methodology Overview}

Evaluating a system whose central purpose is the faithful expression of creative intent is inherently multidimensional. No single metric can capture whether a generated 3D object truly reflects the structural requirements, stylistic inspirations, and conceptual nuances encoded in a moodboard. Accordingly, this evaluation adopts a hybrid strategy: quantitative objective metrics provide directional, reproducible measurements of semantic alignment, while subjective human studies serve as the primary arbiter of preference and perceived quality.

This dual design also provides a direct safeguard against the well-known risk of a system being inadvertently optimised toward its own evaluation metric. Four structural properties of this setup prevent that failure mode. First, the evaluation space is frozen: the embedding model used for all objective metrics is an off-the-shelf, pretrained CLIP model that is never fine-tuned on any outputs from this system. Second, comparisons are relative rather than absolute: the exact same metric is applied identically to both Imagin3D and the baseline, so any systematic bias in the metric affects both arms equally and does not distort the comparison. Third, cross-metric validation is built in: because the two objective metrics measure orthogonal properties (2D-to-3D fidelity versus moodboard semantic alignment) and both are cross-checked against the subjective study, a spuriously high score on one metric cannot mask a poor result on another. Fourth, and most critically, the final judgment is delegated to human participants: the objective scores are treated as directional proxies rather than ground truth, and the subjective A/B preference test determines whether participants actually perceive a meaningful quality difference.

\section{Baseline Workflow}
\label{sec:baseline}

\subsection{The Cascaded-Design Constraint}

The state-of-the-art in direct 3D generation is uniformly shaped by a single architectural constraint: regardless of the model family, all current high-quality approaches impose a 2D image bottleneck at some stage of the pipeline. This is not an incidental limitation but a structural one. Systems such as Trellis v2, Hunyuan3D-2, and the broader family of image-conditioned latent diffusion pipelines require a single high-fidelity reference image as their primary conditioning signal. While some variants of these architectures nominally accept a text prompt through an online interface, all of them still internally resolve that text to a single intermediate image before 3D generation begins — a design pattern appropriately described as cascaded. No current production-ready model accepts a heterogeneous collection of images, 3D models, videos, text, and colour palettes as joint direct inputs for 3D generation without first reducing them to this 2D bottleneck.

This cascaded architecture creates a specific usability problem for designers working from moodboards. Because modern 3D generators accept only a single image, a user wishing to leverage diverse visual references must manually bridge the modality gap themselves: they must translate the collective intent of their reference collection into a text description that an intermediate text-to-image model can process, generate a single image from that description, and only then feed that image into the 3D pipeline.

\subsection{Manual Prompt, Text-to-Image, and Trellis v2}

This manual bridging process defines the experimental baseline. In the baseline workflow, the researcher hand-writes a comprehensive text prompt that distils the full design intent of the moodboard --- capturing the object category, geometric structure, material properties, stylistic characteristics, and colour relationships --- into a single natural-language description. This prompt is then passed to Gemini 2.5 Flash Image (Nano Banana), a state-of-the-art text-to-image model, which generates a 2D reference image. Finally, that image is fed into Trellis v2, the same image-to-3D latent diffusion model used by Imagin3D, to produce the final 3D mesh.

An important practical limitation of this baseline also deserves explicit mention. Because the leading 3D generation models do not yet reliably accept arbitrary video, 3D models, or complex multi-image layouts as direct inputs, a user who wishes to reference these richer modalities has no automated path for doing so. They must manually extract relevant information --- for instance, by taking screenshots of a 3D model from multiple angles, or by identifying and capturing representative keyframes from a video --- before even beginning the text-prompt authoring process. This overhead is not captured in the baseline's text prompt itself, but it represents a real hidden cost in the manual workflow that would compound with the complexity of the moodboard.

\subsection{Fair Parity}

To isolate the contribution of Imagin3D's orchestration pipeline as precisely as possible, the two arms of the comparison are made maximally equivalent in all respects other than the conditioning mechanism. Both arms use the same text-to-image model (Gemini 2.5 Flash Image) and the same 3D generator (Trellis v2). Both are run with identical generation seeds to reduce stochastic variance. The sole experimental variable is therefore whether the input image fed to Trellis v2 was produced by Imagin3D's automated orchestration pipeline --- which performs spatial weighting, intent-based routing, clustering, master-prompt synthesis, and style-image injection --- or by the manually authored text prompt alone. This ablation directly tests the central thesis: that structured moodboard orchestration produces superior 3D conditioning signals compared to human-engineered manual translation.

\section{Test Dataset and Fixed Pipeline}

\subsection{Moodboard Dataset Composition}

To ensure a comprehensive evaluation, the test dataset consists of a curated collection of moodboards spanning a range of object categories and levels of stylistic complexity. The dataset covers both structurally rigid objects --- such as furniture and architectural elements with well-defined geometric constraints --- and organic or highly stylised subjects, including characters and abstract design concepts. This range is deliberate: rigid objects provide a cleaner signal for evaluating geometric fidelity, while stylised subjects stress-test the system's ability to harmonise abstract visual metaphors and non-standard design vocabularies.

Each moodboard in the dataset was authored in the Imagin3D interface and saved as a structured JSON file alongside the raw asset files it references. The JSON schema captures each element's type (image, video, 3D model, text, or colour palette), canvas position, bounding-box scale, explicit cluster membership, and the manually authored baseline text prompt. Because the baseline prompt is hand-written once per moodboard rather than automatically derived, it represents a best-effort human-level translation of the moodboard's design intent into a single text description.

\subsection{Fixed Pipeline Specification}
\label{sec:fixedpipeline}

Both arms of the evaluation are driven by a single headless evaluation script (\texttt{pipeline/run\_ab.py}) that takes a moodboard JSON file as input and produces a pair of 3D GLB meshes --- one from Imagin3D and one from the baseline --- as output. Crucially, both arms consume the exact same raw assets. For Imagin3D, these assets are processed by the full orchestration pipeline as described in Chapter~\ref{chap:metho}. For the baseline, only the manually authored text prompt is used; the visual assets on the moodboard are not consulted. The two arms are executed sequentially on a single RTX 4090 GPU, with GPU memory explicitly cleared between runs to prevent residual activations from influencing subsequent generation.

The pipeline archives all intermediate outputs --- the synthesised master prompt, the master image, the GLB mesh, and its Blender-rendered views --- alongside a run manifest that links both arms to their shared moodboard context. These archived outputs serve as the input to the objective metrics, and the paired GLB files are loaded directly into the A/B preference study viewer.

\subsection{Reproducibility}

Stochastic variance in diffusion-based generation can produce meaningfully different outputs from identical prompts. To control this, the text-to-image generation seed is fixed to 42 across all runs for both arms, and the Trellis v2 pipeline seed is fixed to 0. Because the intent-router weighting in the headless pipeline is also deterministic given fixed model weights, a run is fully reproducible from the moodboard JSON file and the fixed seeds. Pipeline code, model checkpoints, and moodboard files are version-controlled together so that any result reported in Chapter~\ref{chap:results} can be independently reproduced.

\section{Measuring Performance}

Evaluating the effectiveness of moodboard-based conditioning presents unique challenges. Understanding spatial relationships and volumetric properties requires mental rotation and multi-perspective evaluation that 2D images do not support. Consequently, the hybrid evaluation strategy described above employs two distinct subjective protocols and two automated objective metrics to obtain a rounded assessment of both output quality and alignment with design intent.

\section{Subjective Metrics}

Because evaluating multi-faceted stylistic guidance is inherently qualitative, human evaluation remains the most critical performance indicator. Two complementary paradigms are used.

\subsection{A/B Preference Testing}

In the A/B preference test, participants are presented with the source moodboard alongside two anonymised, side-by-side 3D models rendered in an interactive WebGL viewer. The left-right assignment of the two models is randomised per case so that neither arm systematically occupies the more visually prominent position. Participants can freely orbit each model before selecting which one better harmonises the disparate influences and overall design intent of the moodboard. They are also invited to leave a brief written justification for their choice.

The study is delivered through a dedicated browser interface accessible at \texttt{/AB} on the same server that hosts the Imagin3D frontend. This co-location ensures that the same participants who evaluate A/B preference can also directly compare the interactive moodboard tool with the manual baseline workflow, allowing preference ratings to be contextualised within a broader assessment of the overall system.

\subsection{Targeted Recreation Task}

To evaluate the system's utility in professional creative pipelines, a complementary study asks participants to recreate a specific target visual concept using both the manual baseline workflow and the Imagin3D moodboard interface. This task is evaluated across two dimensions. The first is time-to-completion: the total time from starting the task to the moment the participant judges their result satisfactory, measuring whether the moodboard paradigm reduces the traditional generate-and-hope loop inherent to current 3D workflows. The second is intuitiveness, rated on a five-point Likert scale capturing how naturally participants could express their structural and stylistic intent using each approach.

\section{Objective Metrics}

\subsection{The Reconstructive Nature of 3D Generative Models}

A prerequisite for interpreting the objective metrics correctly is an understanding of a fundamental asymmetry between 2D and 3D generative models. Whereas modern 2D diffusion models have been trained on enormous, semantically diverse image datasets and have developed strong priors for creative interpolation, current 3D generative models are primarily trained to reconstruct geometry from images. This reconstructive orientation means that 3D models are generally more conservative than their 2D counterparts: they tend to reproduce the explicit visual content of their conditioning image rather than creatively synthesise novel forms from abstract stylistic cues. This limitation is not a failure of the models themselves but a consequence of the training data available for 3D generation, which skews heavily toward photographed objects rather than stylised or abstract designs.

The practical implication for this evaluation is that the quality ceiling for any metric measuring alignment between a moodboard and a generated 3D object is lower than it would be for 2D generation. The 3D model can only express design intent to the degree that the intermediate 2D image captures it first; any ambiguity or abstraction not resolved at the 2D stage will propagate unresolved into the 3D output. The objective metrics reported below should therefore be understood as measuring the quality of the 2D-to-3D pipeline given the conditioning signal it received, rather than the expressive range of the 3D model itself.

\subsection{Embedding Choice: CLIP over ULIP-2}

Measuring semantic alignment between unstructured 2D and textual inputs and 3D outputs requires mapping both into a shared embedding space. An initial candidate for this was ULIP-2, which extends contrastive pre-training across text, image, and 3D point cloud modalities. In principle, ULIP-2's trimodal alignment would allow direct comparison of a 3D mesh against image and text references in a single unified space. In practice, however, ULIP-2 suffers from a fundamental generalisation problem: the difficulty of aligning more modalities grows with the number of modalities, and the 3D branch of ULIP-2 is trained on Objaverse, a dataset that is heavily biased toward clean, isolated CAD-like objects. When confronted with niche design styles, stylised forms, or the abstract visual metaphors typical of creative moodboards, the 3D embedding space becomes brittle and produces unreliable similarity estimates.

CLIP, by contrast, has been pre-trained on hundreds of millions of image-text pairs drawn from the open web, giving it substantially stronger generalisation to novel concepts, non-standard aesthetics, and diverse stylistic vocabularies. While CLIP does not directly embed 3D geometry, this limitation is readily addressed by rendering the generated 3D model from multiple viewpoints with Blender and embedding the resulting images. This render-then-embed strategy converts the 3D evaluation problem into a 2D image comparison problem in which CLIP's generalisation strengths are directly applicable. The resulting metrics are more reliable on the kinds of stylised, moodboard-driven outputs that constitute this evaluation's test cases.

All objective metrics in this work use the CLIP ViT-L/14 model pre-trained on the OpenAI dataset, as provided by the OpenCLIP library. This model is never fine-tuned on any data from this system; its weights are frozen throughout all experiments.

\subsection{2D-to-3D Preservation}

The first objective metric measures the fidelity of the 3D uplift: how faithfully Trellis v2 reproduced the 2D master image as a 3D object. The generated GLB mesh is rendered from three orbital viewpoints using Blender, and the CLIP image embeddings of these renders are averaged into a single model embedding. This model embedding is then compared against the CLIP image embedding of the master image using cosine similarity. A high score indicates that the 3D generator successfully captured the geometry, material, and overall appearance of its 2D conditioning image. A low score suggests that the model hallucinated or significantly reinterpreted the input, or that the master image contained ambiguities that led to a structurally inconsistent 3D output.

This metric is computed identically for both arms. For Imagin3D, the master image is the output of the Visualizer agent. For the baseline, the master image is the Nano Banana image generated from the manually authored text prompt. Comparing the preservation scores across both arms therefore also provides an indirect signal about which arm produced a cleaner, more 3D-compatible conditioning image.

\subsection{Moodboard Closeness}

The second objective metric measures semantic alignment between the generated 3D model and the original moodboard. For each moodboard element that received a relevance weight above the routing threshold of 50 from the intent router, a CLIP embedding is computed using the modality-appropriate encoder: image and rendered model elements are encoded with the CLIP image encoder, while text and colour palette elements are encoded with the CLIP text encoder and an image encoder applied to a rendered swatch strip, respectively. A weighted centroid of these element embeddings is then computed using the intent-router weights as the weighting function, yielding a single vector that represents the collective design intent of the moodboard as perceived by CLIP.

The closeness score is the cosine similarity between this weighted centroid and the CLIP image embedding of the generated model's renders. This metric directly answers the question of whether the final 3D output has moved toward the design intent encoded in the moodboard. Crucially, both arms are scored against the same centroid, computed from the same elements with the same weights. This allows a direct comparison: Imagin3D's generated mesh versus the baseline's generated mesh, measured against a shared, moodboard-derived semantic target. The difference in closeness scores between the two arms provides a quantitative signal for whether Imagin3D's orchestration pipeline produces 3D outputs that are semantically closer to the moodboard's collective intent than the manually prompted baseline.

\subsection{Safeguards Against Metric Overfitting}

As noted in Section~\ref{sec:baseline}, neither the embedding model nor any component of this evaluation was tuned or selected based on outputs from this system. The evaluation methodology was fixed prior to generating experimental results, and the metrics are applied with identical hyperparameters to both arms. Any systematic bias in the CLIP embedding space --- for instance, a tendency to overestimate similarity for certain object categories --- affects both arms equally and therefore does not distort the relative comparison. The presence of two structurally independent objective metrics provides a further check: a result that is directionally consistent across both preservation and closeness is substantially more credible than one that appears in only one metric. Finally, and most importantly, the objective metrics are explicitly positioned as supporting evidence rather than the primary outcome. The A/B subjective preference test, in which human participants make holistic evaluative judgements with no access to the metric values, constitutes the definitive evaluation of whether Imagin3D produces perceptibly better 3D outputs from a given moodboard.
